\section{Related Work}
\label{sec:related}

The question of when a pure program may update in place is as old as functional arrays. We organise the
comparison by the two choices that define our design: the property is \emph{uniqueness at the site,
decided by liveness}, and a failure is \emph{a compile-time error, not a fallback}.

\paragraph{Linearity and uniqueness.} Wadler showed that a linearly used array can be updated in place,
and introduced \code{let!} to read a linear value for a while without consuming it
\citep{wadler1990linear}; a year later he observed that linearity alone does not imply a single
reference, and stated the goal that ``only update operations are sequentialised but lookups can occur in
parallel'' \citep{wadler1991use}. Marshall, Vollmer and Orchard make the distinction precise: linearity
restricts what can be done with a value in the future, uniqueness guarantees what has been done with it
in the past \citep{marshall2022entente}. Our property is a uniqueness property at one program point,
and our ``read before the write'' plays the part of \code{let!} with no construct, because strict
evaluation order makes the reads' position syntactic. Boyland's alias burying lets a unique variable be
read without being destroyed provided that every alias is dead before the unique variable is used again,
checked by a liveness-style analysis rather than by types \citep{boyland2001alias}; it is the closest
ancestor of our definition. Linear Haskell attaches linearity to function arrows and proves that its
mutable-array API is observationally equivalent to a pure semantics \citep{bernardy2018linear}; the user
writes the linear arrows and the array API's continuation style.

\paragraph{Clean.} Clean's uniqueness typing puts the attribute \code{*} on types, with a subtyping
relation between unique and non-unique types, propagation through containers, and an occurrence-based
marking of variables refined by observer references for guards and strict \code{let}s
\citep{smetsers1994guaranteeing,barendsen1993conventional,barendsen1996uniqueness,plasmeijer2011clean}.
Barendsen and Smetsers show that the system has no principal uniqueness types and treat an expression
whose typing would need a choice among instances as untypable \citep[\S8]{barendsen1996uniqueness}; de
Vries et al.\ reformulate the system to fit standard inference, and still require every binding in a
recursive group to be non-unique \citep{devries2007redefined,devries2008simplified}. They call partial
application ``probably the most subtle aspect of uniqueness typing'' \citep[\S4.2]{devries2008simplified}.
We take from Clean the placement of the demand on library primitives and the notion of deep uniqueness;
we differ in three ways. Our facts are flows, not attributes on types, so there is no principal-type
problem and recursive groups get their best summaries by a fixpoint. Our uniqueness is decided by
liveness rather than by occurrence counts, so reads before a write and dead aliases are free. And beni has
no partial application to make subtle.

\paragraph{Futhark.} Futhark is the closest production design in intent: a pure array language that
guarantees in-place updates statically and reports a violation as an error
\citep{henriksen2017futhark,henriksen2017thesis}. Its consumption rule for sequenced expressions
requires that what the first consumes be neither observed nor consumed by the second,
$(\mathcal{O}_2 \cup \mathcal{C}_2) \cap \mathcal{C}_1 = \emptyset$
\citep[Fig.~5, \textsc{Occurence-Seq}]{henriksen2017futhark}, which is the one-line form of our \Dead\
condition under syntactic order. Futhark differs from our design in requiring uniqueness annotations
(\code{*}) on consuming parameters and unique results, and in an intraprocedural alias analysis, for which
``we are forced to be conservative'' \citep[\S3.2]{henriksen2017futhark}. Henriksen's 2026 retrospective
argues that other language designers should avoid making a type system reason about aliasing, and
recommends reference counting with copy-when-shared for most languages \citep{henriksen2026aliasing}; the
same year they moved alias checking out of the type checker into a pass over fully typed programs
\citep{henriksen2026rewrite}. Our design is such a pass, with interprocedural summaries, and is not part of
the type system.

\paragraph{Usage aspects and destructive-update analysis.} Aspinall, Hofmann and Kone\v{c}n\'y's usage
aspects---destroyed, shared with the result, read only---are our $\mathit{use}$ lattice, and their
soundness theorem for in-place evaluation against a safe semantics is the shape of
Theorem~\ref{thm:main} \citep{aspinall2002another,aspinall2008usage}. Their system is first-order,
monomorphic, and uses explicit resources for allocation; higher-order functions are deferred to a types
and effects extension, which our ownership classes are. Wand and Clinger optimise destructive array updates
with set constraints: an update is made destructive when the array ``can't be bound to a live location''
after its arguments are evaluated, and the constraints ``can be solved in polynomial time''
\citep{wand2001set}. Their analysis is a first-order optimisation that falls back to a copy. Sisal's
compiler removed copies by analysis and reference counts \citep{cann1990sisal}, as does SAC
\citep{grelck2006sac,grelck2004implicit}; both are optimisations with a run-time fallback. Mercury's unique
modes let a programmer declare when a value has one reference and when it will have no more
\citep{mercury2026unique}.

\paragraph{Reference counting with reuse: Lean, Koka, Roc and Morphic.} Lean~4 updates arrays in place
when their reference count is one, and infers which parameters may be borrowed (passed without a count
increment) by a fixpoint that starts from ``all borrowed'' \citep{ullrich2019beans,demoura2021lean4}.
Perceus gives Koka precise reference counting with reuse \citep{reinking2021perceus}, refined by
frame-limited reuse \citep{lorenzen2022frame}. FP\textsuperscript{2} adds a static check, \code{fip}, that
a function is fully in place \citep{lorenzen2023fp2}; but ``the fip annotation in Koka only guarantees that
no (de)allocation occurs if the parameters are unique at runtime'' \citep[\S3.1]{lorenzen2024bst}, so the
guarantee is conditional on a fact checked only by the count. Roc's compiler uses exact reference counts
to mutate in place \citep{teeuwissen2023roc}. Brandon et al.'s Morphic infers borrows with lifetimes with no
annotations and, when a program cannot be typed, ``inserts a handful of reference count operations so that
it can be typed'' \citep{brandon2026borrows}. All of these accept every program and move the cost into
run-time counts or copies. Huisinga's static uniqueness analysis for Lean is the closest to an
error-producing design in this family; it treats higher-order functions as always shared, which he calls a
considerable limitation \citep[\S4.1]{huisinga2023lean}. In a garbage-collected JavaScript target, decrements
buy little: our prototype study found that a sticky ``shared'' bit with no decrements made the same in-place
decisions as a full count in every scenario, and that neither could help a UI runtime that retains what it rendered.

\paragraph{OxCaml modes.} OxCaml adds modes---locality, uniqueness, linearity and more---to OCaml, with
\code{once} closures for those that consume a unique capture, a lock rule for closures, and inference within
implementations \citep{lorenzen2024oxidizing,oxcaml2025uniqueness}. Mode crossing exempts types whose values
cannot interact with an axis, without which adoption ``would have been infeasible'' in a large codebase
\citep[\S1]{peters2026mode}; our crossing rule is the same idea computed from beni's solved types. Two
differences matter. Modes are written in interfaces, while our summaries are inferred and published. And
OxCaml's authors note that the complications of currying for uniqueness ``go away if we imagine a change to
the language forbidding currying'' \citep[\S6.4]{lorenzen2024oxidizing}---which is beni.

\paragraph{Rust.} Rust's borrow checker is the most widely used ownership system, and two of its refinements
are ours in another setting. Non-lexical lifetimes base borrows on liveness \citep{rfc2094}, and two-phase
borrows let a mutable borrow act as a shared one until it is activated \citep{rfc2025}; Polonius recasts
the analysis as Datalog over directional facts \citep{stjerna2020polonius}. Formal models of the borrow
checker exist at several levels \citep{weiss2019oxide,pearce2021lightweight}. Rust has references into
structures, lifetimes in signatures and explicit annotations; beni has none of these, so the problem cases
that motivate much of Rust's machinery (a reference to a slot held across a conditional) do not arise, and
the remaining fact---who may still read a list---can be inferred. Emre et al.'s measurement of
unification-based versus subset-based alias analysis in the translation of C to Rust is the reason our
analysis is directional \citep{emre2023aliasing}. Crichton's studies of how programmers understand ownership
errors shaped our messages \citep{crichton2020usability,crichton2023grounded,crichton2024profiling}.

\paragraph{Swift and Hylo.} Swift's value types are copy-on-write, testing uniqueness at run time; its
ownership manifesto notes that ``reference counting and uniqueness testing do impose some overhead'' and
that performance ``can still be complex to analyze and predict'' \citep{mccall2017manifesto}. Later
proposals added explicit \code{borrowing} and \code{consuming} parameter modifiers and a \code{consume}
operator \citep{swift2022se0377,swift2022se0366}, and a 2025 discussion of an explicit-copies mode set out
the stability requirement we adopt \citep{swift2025explicitcopies}. Mutable value semantics, in which a law of exclusivity makes it safe to sidestep conceptual copies, has been
studied as an implementation strategy \citep{racordon2022mvs}, and Hylo is a language built on it; in Hylo ``copies are explicit by
default'', and the compiler offers to insert them \citep{hylo2024bindings}. Our error and its \code{List.copy}
fix have the same spirit in a pure language with no parameter conventions.

\paragraph{Reachability types and capture tracking.} Reachability types track, in a function's type, which
arguments its result may reach, with a freshness marker for new values
\citep{bao2021reachability,wei2024polymorphic}; recent work adds bidirectional inference given annotations
on function arguments \citep{jia2026escape}, flow-sensitive use and kill effects for ownership transfer through
closures \citep{deng2025capture}, and separation checking that rejects only interfering uses
\citep{xu2026capybara}. O'Connor et al.\ state the guarantee of uniqueness as separation
\citep{oconnor2026separation}. Our $\mathit{res}$ component is a reachability qualifier on results, our
$\mathit{cap}$ facts are the use and kill effects of a closure, and $\fresh$ plays the freshness marker's part.
We infer these facts where those systems declare them, which is possible because beni's lambdas are visible
and its functions first-order in their arrows.

\paragraph{Usage and escape analysis.} Usage analyses for lazy languages infer how often a value is used
\citep{turner1995once,wansbrough1999once,wansbrough2002thesis}; Wansbrough and Peyton Jones's ``poisoning
problem''---one call of a function poisoning the usage of all others---is why our summaries are polymorphic
over function-typed parameters. Blanchet's escape analysis distinguishes closures that are applied from
closures that escape \citep{blanchet1998escape}, which is our $\mathit{escapes}$ bit. Radanne et al.'s Affe
captures borrows in closures with kinds \citep{radanne2020kindly}.

\paragraph{Ownership inference in object-oriented languages.} Clarke et al.'s survey observes that
ownership-type inference approaches ``infer any solution that satisfies the constraints, but cannot give a
best solution'' \citep[\S5]{clarke2013survey}. Our constraints are monotone inclusions over a finite lattice
with a least solution, which the optimistic start finds.

\paragraph{Transient and copy-on-write structures.} Clojure's transients let code build a persistent
structure through a mutable view in $O(1)$ conversion each way, with operations that ``are not designed to be
bashed in-place'' and must thread their return values \citep{hickey2009transients}---a run-time discipline for
the building pattern our checker proves statically.
